\chapter*{CONCLUSION GÉNÉRALE}
\addcontentsline{toc}{chapter}{CONCLUSION GÉNÉRALE}

Ce mémoire visait à démontrer l’intérêt de l’agriculture de précision pour optimiser l’irrigation et la fertilisation en sols salins, à travers une étude comparative menée sur le champ d’expérimentation Agro TIC de l’USSEIN à Kaolack. Une zone test, conduite selon une approche intégrée (étude du sol, amendement ciblé, irrigation et fertilisation pilotées par capteur), a été comparée à une zone témoin conduite selon les pratiques traditionnelles, sur le navet, le gombo et la laitue.

Les plots tests ont conservé un sol meuble et perméable, sans efflorescences salines, avec une meilleure croissance (gombo : 41,3 cm contre 21 cm) et une production effective (navet : 4,8 kg sur 2 m², récolte de laitue), alors que les plots témoins ont donné une production très limitée, voire nulle. Ces résultats vont dans le sens de l’hypothèse H1. Le pilotage par capteur a aussi réduit la consommation d’eau de 66~\%, 76~\% et 86~\% selon les cultures, ce qui confirme l’hypothèse H2. 
Le suivi des paramètres tels que le $pH$, l'humidité et la température du sol et/ou de l'air a également permis de recueillir des informations utiles à la détection précoce d'éventuelles maladies. Ces résultats permettent de confirmer l'hypothèse H3. L'hypothèse H4 n’a pas pu être vérifiée. L’application AgriSens automatise les calculs d’évapotranspiration et de besoins en irrigation et en fertilisation.

Cette étude présente toutefois des limites : le capteur 8-en-1 estime le $NPK$ à partir de l’$EC$, qui n’indique pas la nature des ions présents ; les traitements contre les maladies se sont révélés peu efficaces, et leur détection reposait uniquement sur l’observation.

Ces limites ouvrent des perspectives telles que : utiliser un capteur donnant le $NPK$ réel, identifier précisément les ions du sol, rechercher des remèdes plus efficaces adaptés au contexte local, confirmer les maladies en laboratoire, quantifier les intrants et les coûts, tester l’adoption de l’application par les agriculteurs, et étendre l’étude à d’autres cultures.